\documentclass[11pt,a4paper]{article}

\usepackage[utf8]{inputenc}
\usepackage[T1]{fontenc}
\usepackage[english]{babel}
\usepackage[top=2.2cm,bottom=2.2cm,left=2.2cm,right=2.2cm]{geometry}
\usepackage{booktabs}
\usepackage{tabularx}
\usepackage{array}
\usepackage{ragged2e}
\usepackage{enumitem}
\usepackage[hidelinks]{hyperref}
\usepackage{url}

\newcolumntype{Y}{>{\RaggedRight\arraybackslash}X}
\newcolumntype{L}[1]{>{\RaggedRight\arraybackslash}p{#1}}

\setlength{\parindent}{0pt}
\setlength{\parskip}{0.5em}

\title{\vspace{-1.5cm}Dataset comparison for leaf descriptor extraction}
\author{}
\date{}

\begin{document}
\maketitle
\vspace{-1.5cm}

\section*{Introduction}

When measuring leaf traits -such as aspect ratio, margin serration, color and venation-
there are some relevant features from a dataset that may make the descriptors extraction easier.
The criterion followed in this project compares ten different characteristics among the Flavia,
Pl@ntNet300k, PlanCLEF and Herbarium 2021 datasets.

These ten features are the number of images each dataset contains, the type of the images,
what species have been included, the capture conditions, leaf availability, background, scale,
resolution, annotations and licence/access.

\section*{Comparison}

\begingroup
\scriptsize
\renewcommand{\arraystretch}{1.25}
\setlength{\tabcolsep}{3.5pt}

\begin{tabularx}{\linewidth}{L{2.05cm} Y Y Y Y}
\toprule
 & \textbf{Flavia} & \textbf{Pl@ntNet-300K} &
\textbf{PlantCLEF (2015--16 edition)} & \textbf{Herbarium 2021} \\
\midrule

\textbf{Images} &
1,907 &
306,146 (244k/31k/31k split) &
113,205; other editions range from $\approx$1.4 M to $\approx$4 M &
2,500,779 \\
\addlinespace[3pt]

\textbf{Image type} &
Single isolated leaf, scanned &
Field photographs (Pl@ntNet citizen science) &
Field photographs plus leaf scans on a uniform background &
Digitised herbarium sheets \\
\addlinespace[3pt]

\textbf{Species / taxa} &
32 species (Yangtze delta, China) &
1,081 species &
$\approx$1,000 species (France and neighbouring regions) &
64,500 taxa, normalised to LCVP \\
\addlinespace[3pt]

\textbf{Capture conditions} &
Controlled and uniform &
Uncontrolled: variable light, angle, distance; occlusions &
Mixed: uncontrolled for photographs, controlled for the LeafScan subset &
Semi-controlled, but pressed and aged material \\
\addlinespace[3pt]

\textbf{Leaf availability} &
100\%, one leaf per image, no petiole &
Multi-organ; a per-image organ tag allows a leaf subset to be selected &
Multi-organ, with explicit \textit{Leaf} and \textit{LeafScan} view types:
13,367 + 12,605 training images &
Usually present, but overlapping, folded or fragmented \\
\addlinespace[3pt]

\textbf{Background} &
Uniform white &
Natural and arbitrary &
Natural for photographs; uniform for LeafScan &
Uniform pale sheet; labels automatically blurred \\
\addlinespace[3pt]

\textbf{Scale} &
None &
None &
None &
Yes: millimetre ruler and colour chart on the sheet \\
\addlinespace[3pt]

\textbf{Resolution} &
$\approx$1600 $\times$ 1200 px &
570 $\times$ 570 px on average (180--900) &
Variable, typical digital-camera sizes &
$\approx$6000 $\times$ 4000 px originals, released at 1000 px \\
\addlinespace[3pt]

\textbf{Annotations} &
Species label only &
Species (expertise-weighted user votes), organ, author, licence; official
splits &
Species, genus, family, view type, observation ID, user quality vote, date,
locality, author &
Taxon, family, order; accepted / unresolved status \\
\addlinespace[3pt]

\textbf{Licence / access} &
Free download (SourceForge), cite Wu \textit{et al.} (2007) &
CC BY 4.0 (Zenodo) &
Creative Commons; download from the Pl@ntNet / LifeCLEF servers after
registration &
Kaggle (FGVC8), public after registration \\

\bottomrule
\end{tabularx}
\endgroup

\section*{PlantCLEF in detail}

PlantCLEF is not one dataset but a series of yearly editions whose task changes
substantially, so the choice of edition matters more than the name:

\begin{itemize}[leftmargin=1.4em,itemsep=0.2em,topsep=0.2em]
  \item \textbf{2014--2016 (Pl@ntView).} Around 1,000 species from France and
        neighbouring regions, with a view type annotated per image: Branch,
        Entire, Flower, Fruit, Stem, \textit{Leaf} and \textit{LeafScan}. In
        the 2015 release this gives 13,367 leaf photographs and 12,605 leaf
        scans in the training set alone. These are the editions relevant here.
  \item \textbf{2017.} Scaled up with noisy web-crawled images; organ
        annotation becomes less reliable.
  \item \textbf{2020--2021.} Cross-domain task: identify field photographs from
        herbarium sheets. Interesting as a precedent for Herbarium 2021, not as
        a source of leaf images.
  \item \textbf{2022--2023.} Global scale, $\approx$80,000 species, combining
        expert-verified and semi-reviewed web data. Diversity, but no gain in
        measurable leaves.
  \item \textbf{2024--2025.} Multi-label identification in 0.5 m\textsuperscript{2}
        vegetation quadrats, where leaves overlap and cannot be isolated. Not
        usable for measurement, although the DINOv2 model released with the
        2024 edition may be useful as a feature extractor.
\end{itemize}

A reasonable leaf subset can therefore be selected from PlantCLEF, and the
\textit{LeafScan} category is the interesting part: leaves photographed or
scanned against a uniform background, i.e.\ Flavia-like conditions but with
$\approx$1,000 species instead of 32. That gives an intermediate step between
the controlled and the natural setting rather than a jump straight from one to
the other.

\section*{PlantCLEF or Pl@ntNet-300K for the natural-conditions stage?}

Both provide a way to select leaves, and neither provides a scale, so the
decision rests on what each allows us to \emph{evaluate}:

\begin{itemize}[leftmargin=1.4em,itemsep=0.25em,topsep=0.2em]
  \item \textbf{Repeatability.} PlantCLEF groups images by observation, so
        several images of the same individual plant can be compared. That makes
        it possible to measure whether a descriptor returns consistent values
        for the same leaf under different capture conditions --- arguably the
        central question when evaluating descriptor extraction in the field.
        Pl@ntNet-300K does not expose this grouping.
  \item \textbf{Image quality.} PlantCLEF carries a user quality vote per
        image, which allows the evaluation set to be stratified by quality
        instead of being filtered by hand.
  \item \textbf{Controlled tier.} The LeafScan subset has no equivalent in
        Pl@ntNet-300K.
  \item \textbf{Scale and convenience.} Pl@ntNet-300K is larger (306k images,
        1,081 species), has official splits, a single CC BY 4.0 licence and a
        one-step download, which makes it the more practical option and the
        better one if a learned segmentation step has to be trained.
\end{itemize}

The two are complementary rather than exclusive: Pl@ntNet-300K as the volume
source for field images, PlantCLEF 2015--16 for the repeatability and
image-quality analyses, and its LeafScan subset as an intermediate validation
step. Confirming this requires two counts that are not published and that we
should compute when downloading: how many Pl@ntNet-300K images carry the
\texttt{leaf} organ tag and how they are distributed per species, and how many
PlantCLEF observations contain more than one leaf image of the same plant.

\section*{Proposed selection}

\begin{itemize}[leftmargin=1.4em,itemsep=0.3em,topsep=0.2em]
  \item \textbf{Flavia} --- controlled environment in which to develop and
        initially validate the descriptor extraction pipeline. Isolated leaves
        on a white background make segmentation trivial, so each descriptor can
        be checked before any variability is introduced. Its 32 species are too
        few to support general claims, which is why a second stage is needed.
  \item \textbf{Natural conditions} --- Pl@ntNet-300K as the main source, with
        the PlantCLEF 2015--16 leaf and LeafScan subsets used for the
        repeatability and quality-stratified analyses, subject to the counts
        above.
  \item \textbf{Herbarium 2021} --- not used now, but kept on record. It is the
        only dataset with a metric reference, so it remains the natural option
        if absolute size is studied later. The current priority is descriptors
        obtainable directly from leaf geometry and appearance.
\end{itemize}

Two consequences follow. Since none of the datasets in use provides a scale,
the descriptors will be dimensionless or relative rather than absolute. And
since none of the four provides ground-truth morphological annotations --- no
segmentation masks, no contour points --- validation will require manually
annotating a small sample.

\section*{Sources}

\begingroup
\small
\begin{itemize}[leftmargin=1.4em,itemsep=0.15em,topsep=0.2em]
  \item Flavia: \url{https://flavia.sourceforge.net/} ---
        Wu \textit{et al.} (2007), \url{https://arxiv.org/abs/0707.4289}
  \item Pl@ntNet-300K: \url{https://zenodo.org/records/4726653} ---
        \url{https://github.com/plantnet/PlantNet-300K}
  \item PlantCLEF 2015: \url{https://www.imageclef.org/lifeclef/2015/plant} ---
        2016: \url{https://www.imageclef.org/lifeclef/2016/plant}
  \item PlantCLEF 2022: \url{https://www.imageclef.org/PlantCLEF2022} ---
        2024: \url{https://www.imageclef.org/PlantCLEF2024} ---
        2024 pretrained model: \url{https://zenodo.org/records/10848263}
  \item Herbarium 2021:
        \url{https://www.frontiersin.org/journals/plant-science/articles/10.3389/fpls.2021.787127/full}
        --- \url{https://www.kaggle.com/c/herbarium-2021-fgvc8}
\end{itemize}
\endgroup

\end{document}
